% =====================================================================
%  CARNET D'ENTRAINEMENT — FICHE 6 : Liaisons et structures de Lewis
%  THÈME 2 — Méthode des 5 étapes et structure de Lewis
%  Fiche TUTO (à lire avant les exercices)
%  Compilation : pdflatex
% =====================================================================
\documentclass[11pt,a4paper]{article}
\usepackage[utf8]{inputenc}
\usepackage[T1]{fontenc}
\usepackage{lmodern}
\usepackage[french]{babel}
\usepackage{amsmath,amssymb,mathtools}
\usepackage{icomma}
\usepackage[version=4]{mhchem}
\usepackage{siunitx}
\sisetup{output-decimal-marker={,}}
\usepackage{xcolor}
\usepackage{colortbl}
\usepackage{tikz}
\usepackage[most]{tcolorbox}
\usepackage{enumitem}
\usepackage{array}
\usepackage{booktabs}
\usepackage{multicol}
\usepackage{fancyhdr}
\usepackage[hidelinks]{hyperref}
\usetikzlibrary{arrows.meta,positioning,calc,shapes.geometric,shapes.misc,fit,backgrounds}
\usepackage[a4paper,margin=2cm,headheight=26pt,headsep=8pt]{geometry}

% -------- palette --------
\definecolor{primary}{RGB}{146,95,160}
\definecolor{primarydark}{RGB}{92,52,104}
\definecolor{defcol}{RGB}{0,114,178}
\definecolor{thmcol}{RGB}{0,140,120}
\definecolor{excol}{RGB}{0,135,90}
\definecolor{methcol}{RGB}{90,90,100}
\definecolor{attncol}{RGB}{200,40,55}
\definecolor{astucecol}{RGB}{198,93,9}
\definecolor{atomO}{RGB}{220,55,45}
\definecolor{atomN}{RGB}{48,80,200}
\definecolor{atomH}{RGB}{235,235,235}
\definecolor{lonepair}{RGB}{120,94,200}
\definecolor{bondcol}{RGB}{60,60,70}
\definecolor{piecol}{RGB}{198,93,9}

% -------- boîtes --------
\tcbset{boxbase/.style={enhanced,breakable,boxrule=0.9pt,arc=2.5pt,
   left=3mm,right=3mm,top=2.2mm,bottom=2.2mm,
   fonttitle=\bfseries,coltitle=white,toptitle=0.6mm,bottomtitle=0.6mm}}
\newtcolorbox{defbox}[1][]{boxbase,colback=defcol!6,colframe=defcol,
   title={\textsc{À retenir}\ifx\relax#1\relax\relax\else\textmd{ : \itshape #1}\fi}}
\newtcolorbox{methbox}[1][]{boxbase,colback=methcol!8,colframe=methcol,sharp corners,
   title={\textsc{Méthode}\ifx\relax#1\relax\relax\else\textmd{ : \itshape #1}\fi}}
\newtcolorbox{exbox}[1][]{boxbase,colback=excol!6,colframe=excol,
   title={\textsc{Exemple}\ifx\relax#1\relax\relax\else\textmd{ : \itshape #1}\fi}}
\newtcolorbox{attnbox}[1][]{boxbase,colback=attncol!6,colframe=attncol,
   title={\textsc{Piège à éviter}\ifx\relax#1\relax\relax\else\textmd{ : \itshape #1}\fi}}
\newtcolorbox{astucebox}[1][]{boxbase,colback=astucecol!8,colframe=astucecol,
   title={\textsc{Astuce}\ifx\relax#1\relax\relax\else\textmd{ : \itshape #1}\fi}}

% -------- macros des quatre grandeurs --------
\newcommand{\nmax}{n_{\max}(e^-)}
\newcommand{\nv}{n_{v}(e^-)}
\newcommand{\nl}{n_{\ell}(e^-)}
\newcommand{\nnl}{n_{n\ell}(e^-)}

% -------- en-tête --------
\pagestyle{fancy}\fancyhf{}
\fancyhead[L]{\small\textcolor{primarydark}{\textbf{Fiche 6 · Thème 2} — Méthode des 5 étapes et structure de Lewis}}
\fancyhead[R]{\small\textcolor{primarydark}{\textsc{Tuto}}}
\fancyfoot[C]{\small\thepage}
\renewcommand{\headrulewidth}{0.8pt}
\renewcommand{\headrule}{\hbox to\headwidth{\color{primary}\leaders\hrule height \headrulewidth\hfill}}

\setlength{\parindent}{0pt}
\setlength{\parskip}{3pt}

\begin{document}

\begin{tcolorbox}[boxbase,colback=primary!12,colframe=primary,
   title={\Large Thème 2 — Méthode des 5 étapes et structure de Lewis}]
\textbf{Objectif du thème.} Construire la \emph{structure de Lewis} d'une molécule ou d'un ion
sans dessiner à l'aveugle : on \emph{compte} d'abord les électrons avec quatre grandeurs, puis on
en déduit le \textbf{nombre de liaisons} et le \textbf{nombre de doublets non liants}. C'est une
méthode mécanique, rapide et sûre — parfaite pour un concours.
\end{tcolorbox}

\section*{1. Les quatre grandeurs à ne jamais confondre}

Tout repose sur un petit jeu de quatre nombres d'électrons. Les deux premiers se calculent
directement à partir de la formule ; les deux derniers s'en déduisent par soustraction.

\begin{defbox}[les quatre formules]
\begin{enumerate}[leftmargin=1.7em,topsep=2pt,itemsep=4pt]
  \item \textbf{Électrons souhaités} — ce qu'il faudrait pour que \emph{chaque} atome ait son
        octet (ou son duet), pris \emph{séparément} :
        \[ \nmax = 8\,a + 2\,b \]
        où $a$ = nombre d'atomes visant l'\textbf{octet} (tous \emph{sauf} \ce{H} et \ce{He}) et
        $b$ = nombre d'atomes visant le \textbf{duet} (les \ce{H}).
  \item \textbf{Électrons disponibles} — ceux réellement présents dans l'édifice :
        \[ \nv = \Big(\textstyle\sum_i N_v^{(i)}\Big) - q \]
        où $N_v^{(i)}$ = électrons de valence de l'atome $i$ (= son numéro de colonne) et $q$ =
        charge algébrique de l'édifice.
  \item \textbf{Électrons liants} $\Rightarrow$ nombre de liaisons :
        \[ \nl = \nmax - \nv \qquad\Longrightarrow\qquad \text{nb de liaisons} = \frac{\nl}{2} \]
  \item \textbf{Électrons non liants} $\Rightarrow$ nombre de doublets libres :
        \[ \nnl = \nv - \nl \qquad\Longrightarrow\qquad \text{nb de doublets libres} = \frac{\nnl}{2} \]
\end{enumerate}
\end{defbox}

\begin{attnbox}[le signe de la charge $q$ dans $\nv$]
On \textbf{retranche} $q$. Retenez le réflexe par le signe de l'ion :
\begin{itemize}[leftmargin=1.6em,topsep=1pt,itemsep=1pt]
  \item \textbf{anion} (\,$q<0$, ex. \ce{CO3^2-} : $q=-2$\,) : $-q$ est \emph{positif}, on
        \textbf{ajoute} des électrons ($\dots+2$) — l'ion en a \emph{gagné} ;
  \item \textbf{cation} (\,$q>0$, ex. \ce{NH4+} : $q=+1$\,) : $-q$ est \emph{négatif}, on
        \textbf{retranche} des électrons ($\dots-1$) — l'ion en a \emph{perdu}.
\end{itemize}
\end{attnbox}

\begin{astucebox}[lire la valence $N_v$ directement dans le tableau]
\centering
\renewcommand{\arraystretch}{1.2}
\begin{tabular}{c|c c c c c c c}
\toprule
Atome & \ce{H} & \ce{B} & \ce{C},\,\ce{Si} & \ce{N},\,\ce{P} & \ce{O},\,\ce{S} & \ce{F},\,\ce{Cl},\,\ce{Br} & \ce{Ne} \\
$N_v$ & 1 & 3 & 4 & 5 & 6 & 7 & 8 \\
\bottomrule
\end{tabular}
\end{astucebox}

\section*{2. Les cinq étapes, dans l'ordre}

\begin{methbox}[l'enchaînement à appliquer]
\begin{enumerate}[leftmargin=1.9em,topsep=2pt,itemsep=2pt]
  \item Écrire la \textbf{formule brute} et repérer l'\textbf{atome central} (le moins
        électronégatif ; \ce{H} n'est \emph{jamais} central). Compter $a$ et $b$.
  \item Calculer $\nmax = 8a + 2b$.
  \item Calculer $\nv = \sum N_v^{(i)} - q$ (attention au signe de $q$).
  \item Calculer $\nl = \nmax - \nv$, puis $\nl/2 =$ \textbf{nombre de liaisons}. Les placer
        (une double compte pour $2$, une triple pour $3$).
  \item Calculer $\nnl = \nv - \nl$, puis $\nnl/2 =$ \textbf{nombre de doublets libres}. Les
        répartir pour compléter les octets, puis \textbf{vérifier} chaque octet (et duet).
\end{enumerate}
\end{methbox}

\section*{3. Exemple guidé : l'ammoniac \ce{NH3}}

\begin{exbox}[les 5 étapes déroulées sur \ce{NH3}]
\textbf{Étape 1 — formule et atome central.}\quad \ce{NH3}, centre = \ce{N} (l'hydrogène ne peut
pas être central). On a $a=1$ (le seul \ce{N}) et $b=3$ (les trois \ce{H}).

\smallskip
\textbf{Étape 2 — électrons souhaités.}\quad $\nmax = 8\times 1 + 2\times 3 = \mathbf{14}$.

\smallskip
\textbf{Étape 3 — électrons disponibles.}\quad Molécule neutre, $q=0$ :
$\nv = \underbrace{5}_{\ce{N}} + \underbrace{3\times 1}_{3\,\ce{H}} - 0 = \mathbf{8}$.

\smallskip
\textbf{Étape 4 — liaisons.}\quad $\nl = 14 - 8 = \mathbf{6}\;\Rightarrow\;\dfrac{6}{2}=3$
liaisons : les trois liaisons \ce{N-H}.

\smallskip
\textbf{Étape 5 — doublets libres.}\quad $\nnl = 8 - 6 = \mathbf{2}\;\Rightarrow\;\dfrac{2}{2}=1$
doublet non liant, porté par l'azote. \emph{Vérification} : \ce{N} voit $3$ liaisons ($6\,e^-$)
$+\,1$ doublet libre ($2\,e^-$) $= 8\,e^-$, octet respecté ; chaque \ce{H} voit $2\,e^-$, duet
respecté. \checkmark
\end{exbox}

\begin{center}
\begin{tikzpicture}[scale=1.0]
  \coordinate (N) at (0,0);
  \coordinate (Hh) at (0,1.4);
  \coordinate (Hl) at (-1.4,-0.85);
  \coordinate (Hr) at (1.4,-0.85);
  \draw[bondcol,line width=1.4pt] (N)--($(Hh)+(0,-0.3)$);
  \draw[bondcol,line width=1.4pt] (N)--($(Hl)+(0.27,0.16)$);
  \draw[bondcol,line width=1.4pt] (N)--($(Hr)+(-0.27,0.16)$);
  \node[font=\large\bfseries] at (N) {N};
  \node[font=\large\bfseries] at (Hh) {H};
  \node[font=\large\bfseries] at (Hl) {H};
  \node[font=\large\bfseries] at (Hr) {H};
  \fill[lonepair] ($(N)+(168:0.62)$) circle (2.4pt);
  \fill[lonepair] ($(N)+(192:0.62)$) circle (2.4pt);
  \node[lonepair,font=\footnotesize,align=center] at (-2.35,0.5) {doublet\\non liant};
  \draw[lonepair,->,>=Stealth] (-1.8,0.45) -- ($(N)+(180:0.68)$);
  \node[draw=excol,rounded corners=3pt,fill=excol!6,font=\footnotesize,align=left,anchor=west]
    at (2.1,0.1)
    {$\nmax = 14$\\[1pt] $\nv = 8$\\[1pt]
     $\nl = 6 \Rightarrow 3$ liaisons\\[1pt]
     $\nnl = 2 \Rightarrow 1$ doublet libre};
\end{tikzpicture}
\end{center}

\begin{attnbox}[ne pas confondre les quatre grandeurs]
$\nmax$ est un idéal (« ce qu'il faudrait »), $\nv$ est la réalité (« ce qu'on a »). Leur écart
$\nl$ mesure les électrons \emph{mis en commun} (car chaque liaison profite à deux atomes). Ce qui
reste, $\nnl$, ce sont les doublets \emph{libres}. Divisez toujours par $2$ pour passer des
électrons aux \textbf{doublets} (= liaisons ou paires libres).
\end{attnbox}

\begin{astucebox}[contrôle de cohérence]
Après coup, vérifiez que $\nl + \nnl = \nv$ : les électrons disponibles se répartissent
\emph{exactement} entre liaisons et doublets libres. Si ça ne tombe pas juste, une des deux
premières grandeurs est fausse (souvent une erreur de signe sur $q$).
\end{astucebox}

\end{document}
